\lecture{9.6}{Nonparametric Identification with Unmeasured Variables}
\label{lec:9-6}

\subsection{Observed and Unobserved Variables}

Let $V$ denote observed variables and $U$ denote unobserved variables. A causal
graph may be defined over $V\cup U$, but the analyst only observes the marginal
distribution $p(V)$. The identification question is:
\[
  \E\{\Ya{a}\}\quad =\quad ?\quad \text{a function of }p(V).
\]
Nonparametric identification means the answer must hold for every distribution
compatible with the graph, without imposing linearity, normality, or other
parametric restrictions.

\subsection{Adjustment Failure}

The adjustment formula works when there is a measured set $Z$ satisfying the
appropriate graphical criterion:
\[
  \E\{\Ya{a}\}=\E\{\E(Y\given A=a,Z)\}.
\]
If an unmeasured variable lies on an open noncausal path from $A$ to $Y$, no
observed adjustment set may block the path.

\begin{figure}[H]
\centering
\begin{tikzpicture}[
  node distance=1.25cm,
  box/.style={draw, rounded corners=2pt, minimum width=.9cm, minimum height=.7cm},
  arr/.style={-{Latex[length=2.3mm]}, thick}
]
\node[box] (U) {$U$};
\node[box, below left=of U] (A) {$A$};
\node[box, below right=of U] (Y) {$Y$};
\draw[arr] (U) -- (A);
\draw[arr] (U) -- (Y);
\draw[arr] (A) -- (Y);
\end{tikzpicture}
\caption{Unmeasured common cause $U$ creates an open noncausal path between $A$ and $Y$.}
\label{fig:ch9-unmeasured-confounding}
\end{figure}

\subsection{Valid Adjustment Sets}

If $U$ is absent, the parents of $A$ in the causal DAG often provide a valid
adjustment set:
\[
  \Ya{a}\ci A\given \pa(A).
\]
If $U$ is present but an observed set blocks all noncausal paths, the effect is
still identified by adjustment. For example, if $Z$ blocks the back-door path
from $A$ to $Y$, then
\[
  \E\{\Ya{a}\}=\E\{\E(Y\given A=a,Z)\}.
\]

\subsection{Non-Identification}

If no valid adjustment set exists, the effect may or may not be identified by
some other functional. Non-identification means more than failure of a favorite
method. It means there exist two causal models that imply the same observed
distribution $p(V)$ but different values of $\E\{\Ya{a}\}$.

\begin{definition}[Non-identification]
A causal parameter $\psi$ is not identified from $p(V)$ in a model $\M$ if there
exist two distributions $P_1,P_2\in\M$ such that
\[
  p_1(V)=p_2(V)
  \qquad\text{but}\qquad
  \psi(P_1)\ne \psi(P_2).
\]
\end{definition}

\subsection{Linear SEM Counterexamples}

A simple linear structural equation model illustrates the problem. Suppose
\[
  A=\alpha U+\epsilon_A,
\]
\[
  Y=\gamma A+\beta U+\epsilon_Y,
\]
with independent errors. The observed association between $A$ and $Y$ can be
written schematically as
\[
  \rho_{AY}=\gamma+\alpha\beta.
\]
The causal effect is $\gamma$, while $\alpha\beta$ is the spurious association
through the unmeasured common cause. The observed data provide one association
but there are two unknown components. Infinitely many pairs
$(\gamma,\alpha\beta)$ can produce the same $\rho_{AY}$.

For example, if $\rho_{AY}=0.6$, then both of the following are compatible with
the same observed association:
\[
  \gamma=0.6,\ \alpha\beta=0,
  \qquad\text{and}\qquad
  \gamma=0.1,\ \alpha\beta=0.5.
\]
The observed distribution alone cannot choose between them.

\subsection{Using Submodels to Prove Non-Identification}

Non-identification in a simple submodel can be enough to prove
non-identification in a larger model. Let $\M_{\mathrm{sub}}$ be a statistical
model contained in a larger model $\M_{\mathrm{super}}$:
\[
  \M_{\mathrm{sub}}\subset \M_{\mathrm{super}}.
\]
If a causal parameter is identified in the larger model, then it must also be
identified in every smaller submodel:
\[
  \psi \text{ identified in } \M_{\mathrm{super}}
  \Rightarrow
  \psi \text{ identified in } \M_{\mathrm{sub}}.
\]
Taking the contrapositive,
\[
  \psi \text{ not identified in } \M_{\mathrm{sub}}
  \Rightarrow
  \psi \text{ not identified in } \M_{\mathrm{super}}.
\]
This is why a linear structural equation counterexample can prove a
nonparametric non-identification result. If even a restricted linear submodel
contains two distributions with the same observed law and different causal
effects, then the larger nonparametric model also contains those two
distributions.

The reverse implication is false. Identification in a linear Gaussian submodel
does not imply identification in the unrestricted nonparametric model. A linear
model may use algebraic restrictions that are not implied by the graph alone.

\subsection{Why Machine Learning Cannot Solve Non-Identification}

Machine learning can estimate complex functions of the observed distribution
$p(V)$. It cannot recover information that is not contained in $p(V)$. If two
causal models have the same observed distribution but different causal effects,
any algorithm trained only on observed data must return the same output for both
models, while the correct causal answers differ.

\begin{keyidea}
Non-identification is not a failure of modeling flexibility. It is a logical
failure of information: the observed distribution does not contain enough
information to determine the causal parameter under the assumed model.
\end{keyidea}
